\section{Discussion}
\label{sec:discussion}

\subsection{Interpreting the BERT Result}

The $r = 0.984$ correlation substantially exceeds our pre-registered threshold and our pre-experimental confidence estimate (0.72). Three factors contribute:

\textbf{Bimodal oracle structure amplifies correlation.} The BERT oracle assigns $r^* = 4$ or $r^* = 64$ only. erank cleanly separates these tiers --- the problem reduces to a binary structural classification that erank solves deterministically.

\textbf{Mechanism matches theory.} FFN intermediate layers (erank $\sim 704$--$710$, many non-dominated directions) need $r=64$ to cover task-relevant signal. Attention layers --- both output projections and query projections (erank $\sim 556$--$590$, concentrated spectra) --- need only $r=4$. Pre-training geometry drives adaptation need.

\textbf{Sample size caveat.} With $n=5$ oracle measurements spanning both tiers, the near-perfect correlation reflects deterministic bimodal structure, not smooth linear relationship. Full coverage of all 72 layers is necessary to confirm correlation over intermediate erank values.

\subsection{Multi-Family Status and Expectations}

DeBERTa-v3-base and ViT-base oracle sweeps are pending. The consistent layer-type structure in erank maps for all three families, combined with ViT's 2.97$\times$ erank range ratio, motivates cross-family oracle sweeps, though the direction and magnitude of correlation cannot be predicted from erank variation alone. We present current results as preliminary evidence rather than a claim of full hypothesis satisfaction.

\subsection{Limitations}

\textbf{L1: Partial oracle coverage.} 5/72 BERT layers measured (6.9\%). Intermediate erank layers ($\sim 620$--$680$) unmeasured; non-linearities may exist.

\textbf{L2: Marginal vs.\ joint oracle.} The PARA oracle assigns rank per layer independently (all others at $r=8$). Globally optimal joint allocation may differ.

\textbf{L3: Encoder-only models only.} Decoder-only LLMs (LLaMA, Mistral) with causal attention have different SVD properties; generalization is untested.

\textbf{L4: Classification tasks only.} Oracle requires a scalar accuracy metric; generation tasks need different oracle definitions.

\textbf{L5: BERT and DeBERTa erank variation below A1 threshold.} BERT-base and DeBERTa-v3-base show CV $\approx 0.04$, below the pre-registered A1 threshold (CV $> 0.05$). Only ViT-base (CV $\approx 0.12$) meets this threshold. The strong oracle correlation for BERT ($r = 0.984$) is driven by oracle bimodality (rank 4 vs.\ 64) rather than erank spread --- erank need not vary widely to discriminate two structural tiers perfectly, but this means continuous rank discrimination across the full erank range remains untested for NLP encoders. If the oracle is not bimodal for full 72-layer coverage, the modest BERT/DeBERTa erank variation (range ratio 1.34$\times$--1.35$\times$) may limit erank's discriminative power as a continuous predictor.

\textbf{L6: Depth confound in oracle sample.} The 5 measured BERT oracle layers span network depths 0--10, with attention layers sampled from shallower positions (encoder layers 0, 3, 6) and FFN layers from deeper positions (encoder layers 8, 10). Consequently, layer type is partially confounded with depth in the current sample --- whether erank predicts oracle rank independently of depth cannot be assessed from 5 layers alone. Partial correlation analysis controlling for depth is planned as a priority analysis once full-coverage oracle sweeps are completed.

\textbf{L7: Stable rank oracle correlation pending.} Stable rank ($\|W_0\|_F^2 / \|W_0\|_2^2$) maps were computed for all BERT layers but oracle correlation for stable rank is not reported in this paper, as oracle sweeps covered only 5 layers and the focus of the current work is erank as the primary predictor. Stable rank vs.\ erank oracle comparison is deferred to the full-oracle experiment.

\subsection{Broader Impact}

Zero-cost structural rank prediction could reduce calibration overhead for LoRA deployment at scale. We recommend treating erank as a prior over rank allocation rather than a substitute for validation until multi-family and multi-task evidence is established.
